\documentclass[12pt]{article}

\usepackage[a4paper, total={6in, 8in}]{geometry}
\usepackage{textcomp}

\begin{document}
\setlength{\parindent}{0pt}

\def \t {\textrightarrow}
\def \v {\vspace{1em}}
\def \bi {\begin{itemize}}
\def \ei {\end{itemize}}
\def \s[#1] {\section*{#1}}
\def \ss[#1] {\subsection*{#1}}
\def \sss[#1] {\subsubsection*{#1}}

\s[Pascoli - Il Fanciullino]
... godono, piango = climax ascendente \\
Solo/solo è chiasmo, ma anche climax asc. \\
Si uniscono due immagini: il fanciullino che lo è per eta, e il fanciullino che è dentro di noi \\
Il desiderio verso esperienze nuove viene meno, la speranza ha lasciato il posto alla delusione \\
Strillare e guaire \t guaire è onomatopeico e animalesco \\
Il fanciullino si stupisce di tutto \t Marino?? il poeta si deve stupire, e questo stupore si trova nella poesia, associato a un utilizzo di figure retoriche (come iperboli??) \\
Arruginire la voce (sinestesia) = percorrere l'eta biologica e come si evolve la vita in questo modo \\
L'angolo d'anima donde esso risuona \t fino a qua c'è un blocco: noi invecchiamo, ma quella voce dentro l'anima si sente e nell eta matura, durante la quale ci spegniamo perche difendiamo la nostra posizione nel mondo senza ascoltare cio che abbiamo dentro \\
Perorare la nostra causa = linguaggio ciceroniano \\
Il fanciullino risuona dentro l'anima ed è come la coscienza \t fanciullino è la voce della coscienza, che parla sempre ma a volte non ascoltiamo \\
Seneca \\
Tintinnio = onomatopea \\
La smania di crescere fa dimenticare il bagaglio di sensibilita che caratterizza il giovane \\
Gorgheggiare e mormorare = onomatopee, e inoltre specifismo con usignolo \\
L usignolo è stato l epiteto di Edit Piaf \\
Montale il male di vivere

\v

Nava \t l'edizioni che contemplano lo studio delle varianti sono curate da lui \\
Il mondo pascoliano puo essere configurato con:
\bi
    \item il nido insanguinato
    \item la natura
    \item paesaggi campestri
    \item decadentismo e la morte che permea la realta
\ei
Vicinanza con natura si avvicina a Leopardi \\
Pascoli coglie anche i bisogni delle masse \t nella sua parabola universitaria viene incarcerato, e si avvicina allo spirito sorelliano negli anni che culminano con "La grande proletario si è mossa" nel 12 \\
È ancora giovane, e parla di un mondo che rivendica la necessita di uscire dal silenzio \t e questo ha destato supore \\
Questo non gli merito l etichetta di poeta civile, ma da luce a sfaccettatura sulla sua produzione \\
La fratellanza e bonta sono prob. di ispirazione socialista, e prob. partecipa all'internazionale \\
Ha superato il positivismo attraversandolo \\
Subisce un ingiustizia, ma non vendico mai gli assassini del padre \t la sua morte rivendica nell animo giustizia \\
Il suo eros è piu volto al tanatos che all eros stesso \t vuole difendere la partria, il nido = irrisolto sessuale \\
Lui risponde a un bisogno di perfezionismo e specifismo linguistico, di cui parlano COntini in "Specifisimo del linguaggio pascoliano" e Ermengardo?? in ?? \\
Il complesso di Edipo con Telemaco, che è il figlio con figura paterna ingombrante \\
Analisi psicologica porta a dire che non ha vissuto una simmetria nelle figure genitoriali \\
Testo di Gianfranco Contini \t "Varianti ed altra linguistica" è una rivista letteraria, in cui ha anche scritto Contini

\v

I "Poemetti" \t il mondo sta cambiando, e la quotidianita non si puo far \\
La poesia deve difendere l uomo dalla ingiustizie, dai mali della societa industriale, e impegno civile di pascoli si configura in sperimentalismo linguistico, e lessico americaneggiante nel poemetto "Italy", in cui difende l'identita di patria \\
"Poemi conviviali" (1904) \t convivium tradizione dantesca, e ulisse \t necessita di un uomo che va oltre le barriere, e a chiudere la attivita dell autore: gli studi su dante, che compie a Castelvecchio \\
Lo studio di pascoli nei confronti di dante presenta ottica filosofico-erudita, non come Pavese \t l'opera è "Sotto il velame" e "La venerabile visione del 1302" ??
\end{document}
